\documentclass[12pt]{book}
\usepackage[utf8]{inputenc}
\usepackage[T1]{fontenc}
\usepackage{amsmath,amssymb,amsfonts}
\usepackage{mathrsfs}
\usepackage{amsthm}
\usepackage{tikz-cd}

% Calligraphic for categories and sheaves
\newcommand{\cat}[1]{\mathcal{#1}}
\newcommand{\sheaf}[1]{\mathcal{#1}}

% Standard math operators
\DeclareMathOperator{\Hom}{Hom}
\DeclareMathOperator{\Ext}{Ext}
\DeclareMathOperator{\Tor}{Tor}
\DeclareMathOperator{\id}{id}
\DeclareMathOperator{\dimk}{dim_k}
\DeclareMathOperator{\Rfst}{\mathbf{R}f_*}
\DeclareMathOperator{\Lfst}{\mathbf{L}f_*}
\DeclareMathOperator{\RHom}{\mathbf{R}\mathcal{H}\textit{om}}
\DeclareMathOperator{\Res}{Res}
\DeclareMathOperator{\Tr}{Tr}
\DeclareMathOperator{\RGamma}{\mathbf{R}\Gamma}

% Theorem environments
\theoremstyle{plain}
\newtheorem{lemma}{Lemma}[section]
\newtheorem{proposition}{Proposition}[section]
\newtheorem{definition}{Definition}[section]
\newtheorem{corollary}{Corollary}[section]
\newtheorem{remark}{Remark}[section]
\newtheorem{motif}{Motif}[section]

\begin{document}

\setcounter{page}{222}

\subsection{Variation 5}
Combining Variations 1 and 3, we give the following definition.

\begin{definition}
A \textit{sheaf of families of supports} on a topological space \(X\) is a sheaf of sets \(\underline{\varphi}\) such that for every open set \(U\subseteq X\), \(\underline{\varphi}(U)\) is a family of supports on \(U\); and for open inclusion \(V\subseteq U\), the restriction map
\[
\underline{\varphi}(U)\to\underline{\varphi}(V)
\]
sends \(Z\in\underline{\varphi}(U)\) to \(Z\cap V\in\underline{\varphi}(V)\).
\end{definition}

\begin{remark}
\begin{enumerate}
\item If \(\varphi\) is a family of supports on \(X\), define the associated sheaf of families of supports \(\tilde{\varphi}\) as the sheaf associated to the presheaf
\[
U\mapsto \{\,Z\cap U \mid Z\in\varphi\,\}.
\]

\item If \(\underline{\varphi}\) is a sheaf of families of supports on \(X\), define the global family of supports
\[
\Gamma(\underline{\varphi}) = \underline{\varphi}(X).
\]

\item In general, the operations \(\Gamma(\cdot)\) and tilde \(\tilde{\phantom{\varphi}}\) are \textit{not} inverses. For example, on a locally compact Hausdorff space, let \(\varphi\) be the family of compact subsets. Then \(\tilde{\varphi}\) is the maximal sheaf of supports, and \(\Gamma(\tilde{\varphi})\) is maximal; but if \(X\) is non-compact, \(\varphi\ne\Gamma(\tilde{\varphi})\).

\item Let \(f\colon X\to Y\) be a finite-type morphism of locally noetherian preschemes, fix integer \(p\). Define \(\underline{\varphi}(U)\) to be all relatively closed \(Z\subseteq U\) such that for every \(y\in f(U)\), the intersection \(Z\cap X_y\) has codimension \(\ge p\) in the fibre \(X_y\). Then in general \(\Gamma(\tilde{\varphi})\ne\varphi\).

Example: \(Y=\operatorname{Spec}k[x,y]\), \(X=\mathbb{P}_Y^1\), \(p=1\). Take \(Z\subset X\) obtained by blowing up the origin of \(Y\): locally \(Z\in\underline{\varphi}\) away from the origin, but \(Z\notin\Gamma(\underline{\varphi})\).

\item Let \(Z\subseteq X\) be stable under specialization. Define codimension of a point \(x\in X\) as the length of the longest chain of proper specializations
\[
x_0\prec x_1\prec\cdots\prec x_n = x.
\]
Then finite unions of closures of points in \(Z\) form a family of supports \(\varphi\), with associated sheaf \(\tilde{\varphi}\). For locally noetherian \(X\) where every closed irreducible subset has a unique generic point, there is a bijection between specialization-stable subsets \(Z\) and families \(\varphi\) satisfying \(\varphi=\Gamma(\tilde{\varphi})\). In such cases we write
\[
\Gamma_Z,\; H_Z^i,\; \underline{\Gamma}_Z
\]
instead of \(\Gamma_\varphi,H_\varphi^i,\underline{\Gamma}_{\tilde{\varphi}}\).
\end{enumerate}
\end{remark}

Now let \(\underline{\varphi}\) be a sheaf of families of supports on \(X\). Define the sheaf
\[
\underline{\Gamma}_{\underline{\varphi}}(\sheaf{F})
\]
whose sections over open \(U\) are \(\Gamma_{\underline{\varphi}(U)}(U,\sheaf{F}|_U)\).
Define the right derived functors
\[
\mathbf{R}\underline{\Gamma}_{\underline{\varphi}},\quad \underline{H}_{\underline{\varphi}}^i(\sheaf{F}).
\]

\begin{motif}
Repeat Motifs A, C as before.
\end{motif}

\begin{motif}
If \(\underline{\varphi}\) is of \textit{global nature} (\(\underline{\varphi}=\widetilde{\Gamma(\underline{\varphi})}\)), then
\[
\underline{H}_{\underline{\varphi}}^i(\sheaf{F})
= \varinjlim_{Z\in\Gamma(\underline{\varphi})} \underline{H}_Z^i(\sheaf{F}).
\]
\end{motif}

\begin{motif}
There is a spectral sequence
\[
E_2^{pq} = H^p\big(X,\underline{H}_{\underline{\varphi}}^q(\sheaf{F})\big)
\Rightarrow H_{\Gamma(\underline{\varphi})}^{p+q}(X,\sheaf{F}),
\]
or in derived categories:
\[
\RGamma_{\Gamma(\underline{\varphi})}
= \RGamma \circ \mathbf{R}\underline{\Gamma}_{\underline{\varphi}}.
\]
\end{motif}

\setcounter{page}{225}

\subsection{Variation 6}
Combine Variations 2 and 3: let \(Z'\subseteq Z\) be closed subsets. Define
\[
\Gamma_{Z/Z'},\;\mathbf{R}\Gamma_{Z/Z'},\;\underline{H}_{Z/Z'}^i
\]
and repeat Motifs A, B, C.

\subsection{Variation 7}
Combine Variations 1, 2, 3: let \(\underline{\psi}\subseteq\underline{\varphi}\) be sheaves of supports on \(X\). Define
\[
\Gamma_{\underline{\varphi}/\underline{\psi}},\;
\mathbf{R}\Gamma_{\underline{\varphi}/\underline{\psi}},\;
\underline{H}_{\underline{\varphi}/\underline{\psi}}^i
\]
and repeat Motifs A, B, C, D, E.

\subsection{Variation 8}
Define purely punctual local cohomology. Let \(x\in X\). Define \(\Gamma_x(\sheaf{F})\) as the subgroup of the stalk \(\sheaf{F}_x\) consisting of germs represented by sections defined near \(x\) with support contained in \(\overline{\{x\}}\). Define derived functors
\[
\mathbf{R}\Gamma_x,\quad H_x^i(\sheaf{F}).
\]
One has
\[
H_x^i(\sheaf{F})
= \big(\underline{H}_{\overline{\{x\}}}(\sheaf{F})\big)_x,
\]
the stalk at \(x\) of the sheafified local cohomology. Repeat Motifs A, C.

\setcounter{page}{226}

\begin{motif}
Let \(X\) be locally noetherian, every closed irreducible subset having a unique generic point. Let \(Z'\subseteq Z\) be specialization-stable subsets such that every point of \(Z\setminus Z'\) is \textit{maximal} in \(X\) (no non-trivial specialization within \(Z\setminus Z'\)). Then for any \(F^\bullet\in D^+(\cat{Ab}(X))\) there is a canonical isomorphism
\[
\underline{H}_{Z/Z'}^\bullet(F^\bullet)
\cong \bigoplus_{x\in Z\setminus Z'} i_x\big(H_x^\bullet(F^\bullet)\big),
\]
where \(i_x(G)\) denotes the constant sheaf \(G\) supported on \(\overline{\{x\}}\).
\end{motif}

\begin{proof}
It suffices to prove for a single sheaf \(\sheaf{F}\in\cat{Ab}(X)\):
\[
\underline{\Gamma}_{Z/Z'}(\sheaf{F})
\cong \bigoplus_{x\in Z\setminus Z'} i_x\big(\Gamma_x(\sheaf{F})\big).
\]
For open \(U\), map a section in \(\Gamma_{Z\cap U}(U,\sheaf{F}|_U)\) to its germs at maximal points \(x\in (Z\setminus Z')\cap U\). By local noetherianity, support of any section is a finite union of closures of points; maximality ensures only finitely many non-zero germs. The map is injective and surjective by germ-section correspondence.
\end{proof}

\setcounter{page}{227}

\subsection{Coda: Spectral Sequence of a Filtered Space}
\begin{motif}
Let
\[
\underline{\varphi}^0 \supseteq \underline{\varphi}^1 \supseteq \underline{\varphi}^2 \supseteq\cdots
\]
be a filtration of \(X\) by sheaves of supports, and let \(F^\bullet\in D^+(\cat{Ab}(X))\) be bounded below. There exists a biregular spectral sequence
\[
E_1^{pq}
= H_{\underline{\varphi}^p/\underline{\varphi}^{p+q}}^\bullet(F^\bullet)
\Rightarrow H^n(X,F^\bullet).
\]
If the filtration is finite (\(\underline{\varphi}^n=\emptyset\) for large \(n\)), the spectral sequence converges strongly.

In derived categories, one has an iterated distinguished triangle filtration:
\[
\RGamma_{\underline{\varphi}^0}\leftarrow
\RGamma_{\underline{\varphi}^1}\leftarrow
\RGamma_{\underline{\varphi}^2}\leftarrow\cdots\leftarrow
\RGamma_\emptyset = \id.
\]
\end{motif}

\begin{proof}
Take an injective resolution of \(F^\bullet\). The filtered complex
\[
I^\bullet = \underline{\Gamma}_{\underline{\varphi}^0}(I^\bullet)
\supseteq \underline{\Gamma}_{\underline{\varphi}^1}(I^\bullet)
\supseteq\cdots
\supseteq \underline{\Gamma}_{\underline{\varphi}^n}(I^\bullet)=0
\]
yields the standard filtered complex spectral sequence.
\end{proof}

\end{document}